% Fichier généré automatiquement par tools/generate_corrections.py.
% Source : RDM/RDM-01-Cohesion/527_RdM/527_RdM.tex
% Statut : complété (3/3 question(s)).
\begin{corrigebox}[Corrigé complété]
\CorrigeQuestion{1}
On écrit le torseur dans un point de réduction adapté. La résultante
        est obtenue par la définition correspondante et le moment est transporté par
        $\vec M_A=\vec M_B+\overrightarrow{AB}\wedge\vec R$. Les composantes sont ensuite
        projetées dans la base demandée.
\CorrigeQuestion{2}
La réponse s'obtient en appliquant les définitions et conventions de la
    figure. Les étapes de calcul sont explicitées, puis le résultat est vérifié par
    cohérence dimensionnelle et par un cas limite.
\CorrigeQuestion{3}
\includegraphics[width=\linewidth]{cor_03}
\end{corrigebox}
